\begin{lemma}[局所自明・群型付けの窓・群精密化]
\label{lem:mcfg-li-group-window}
$\Sigma$ を空でない有限アルファベット，
$h:\Sigma^*\to M$ を有限モノイド準同型とする．
もし $h(\Sigma^+)\in\mathbf{LI}\vee\mathbf{Grp}$ ならば，
ある $k,\ell\ge0$，有限群 $A$，および
準同型 $g:\Sigma^*\to A$ が存在して
\(h\preceq(h_\Sigma^{k,\ell},g)
 \quad(\Sigma^*\text{ 全体上})\)
が成り立つ．
さらに $h(\Sigma^+)\in\mathbf{LI}$ ならば，
群因子を自明群としてよく，
$h\preceq h_\Sigma^{k,\ell}$ となる．
\end{lemma}
\begin{proof}
有限局所自明半群 $S$，有限群 $A$，
部分半群 $T\subseteq S\times A$，
全射半群準同型 $\pi:T\to h(\Sigma^+)$ を取る．
各 $a\in\Sigma$ について $h(a)$ の逆像
$(s_a,g_a)\in T$ を一つ選ぶ．
これらの積により $\rho:\Sigma^+\to T$ を定めると，
$\pi\rho=h|_{\Sigma^+}$ である．
二つの座標準同型を
$s:\Sigma^+\to S$，$g:\Sigma^*\to A$
と書き，$g(\eps)=1_A$ とする．

正語像 $S_0=s(\Sigma^+)$ は局所自明である．
$q=|S_0|+1$ とおく．
有限局所自明半群に対する標準的な恒等式
$S_0^q=S_0E(S_0)S_0$ および
冪等元 $e,f$ と任意の $u\in S_0$ に対する
$euf=ef$~\cite{Pin1986,kuriyamaCFG} より，
長さ $2q$ 以上の語の $s$-値は
長さ $q$ の接頭辞・接尾辞で決まる．
実際，$w=pxz$，$|p|=|z|=q$ と書き，
$s(p)=aeb$，$s(z)=cfd$ とすれば，
$x$ が空でないとき
\(s(pxz)=ae\,(b\,s(x)\,c)\,fd
       =aefd=s(pz)\)
となる．$x=\eps$ の場合も中央の積を
$bc\in S_0$ と解釈すれば同じ恒等式が使える．
長さが $2q$ 未満の語は
$h_\Sigma^{q,q}$ が個別に区別する．
したがって非空語について窓型が一致すれば
$s$-値が一致し，さらに $g$-値も一致すれば
$\rho$-値，よって $h$-値も一致する．

窓型は空語 $\eps$ と非空語を区別するので，
核の包含は $\Sigma^*$ 全体に拡張される．
$k=\ell=q$ とすれば第一の主張が従う．
$h(\Sigma^+)$ 自体が局所自明ならば，
$S=h(\Sigma^+)$ とし，群因子を自明群に取ればよい．
\end{proof}

\begin{corollary}[多項式特徴データの代数的十分条件]
\label{cor:mcfg-li-group-polydata}
空でない有限 $\Sigma$，ファンアウト上界 $f\ge1$，
規則ランク上界 $r_{\max}\ge1$，有限型付け
$h:\Sigma^*\to M$ を固定する．
$\Cmcf{f}{h}$ の標的を生成する縮約済み良形の
ファンアウト高々 $f$・ランク高々 $r_{\max}$ の
文法 $G$ に対し，
$\bar\theta_G=\max\{1,\theta_G\}$ とする．
ランク制限付き再構成作用素は
個数 $O_{f,r_{\max},h}(\|G\|)$ の特徴標本を持ち，
その総長は
\[
\begin{array}{ll}
 \|\Wit(\widetilde G_0)\|_+
 =O_{f,r_{\max},h,\Sigma}(\|G\|^2\bar\theta_G),
 & h(\Sigma^+)\in\mathbf{LI}\text{ の場合},\\[2pt]
 \|\Wit(\widetilde G_0)\|_+
 =O_{f,r_{\max},h,\Sigma}(\|G\|^3\bar\theta_G),
 & h(\Sigma^+)\in\mathbf{LI}\vee\mathbf{Grp}\text{ の場合}.
\end{array}
\]
\end{corollary}
\begin{proof}
$\mathbf{LI}$ の場合，
補題~\ref{lem:mcfg-li-group-window} により
群因子を自明群として窓型付けに帰着し，
定理~\ref{thm:boundary-polydata} を適用する．
$\mathbf{LI}\vee\mathbf{Grp}$ の場合，
同補題による窓・有限群への精密化と
定理~\ref{thm:window-group-polydata} を適用する．
精密化の諸パラメータは固定した型付けと
アルファベットだけに依存し，$G$ には依存しない．
特徴標本の個数についての上界も両方で成立する．
\end{proof}

\begin{remark}
上の条件は\emph{十分条件}であり，
元の文法の厚みに関する多項式上界が成り立つ型付けの
完全な特徴づけではない．
問題~\ref{op:fanout-two} の高ファンアウト下界は未解決である．
空アルファベットでは語は $\eps$ のみであり自明である．
\end{remark}

